\chapter{Projeto 3: O Labirinto}

\begin{objetivos}
\begin{itemize}
\item Desenhar um cenário com paredes e uma meta;
\item Fazer o personagem andar sem atravessar as paredes;
\item Detectar a vitória e medir o tempo.
\end{itemize}
\end{objetivos}

\section{O plano do labirinto}

Neste jogo o gatinho precisa \textbf{sair de um labirinto} sem encostar nas paredes. Se tocar numa parede, volta ao início. Se chegar à bandeira verde, ganha!

\begin{definicao}[Roteiro do labirinto]
\textbf{Cenário:} paredes \textbf{pretas} e uma meta \textbf{verde}.\\
\textbf{Ator:} Gatinho pequeno (use \bloco[scLooks]{mude o tamanho para \entrada{30} \%}).\\
\textbf{Regras:} tocar na cor preta reinicia; tocar na cor verde vence.
\end{definicao}

\begin{missao}[Missão 19: Escape do labirinto]
Desenhe o labirinto e programe o gatinho para sair dele.
\end{missao}

\section{Passo 1: o cenário}

No editor de desenho do Palco, use o pincel \textbf{preto} bem grosso para desenhar as paredes e pinte um quadrado \textbf{verde} no final do caminho. As cores são importantes, porque o programa vai reconhecê-las!

\begin{dica}
Deixe os corredores \textbf{mais largos} do que o gatinho. Se estiverem muito apertados, o jogo fica quase impossível.
\end{dica}

\section{Passo 2: o movimento}

Use a técnica do capítulo de Animação, mas sem o laço de repetição, apenas um evento para cada seta:

\begin{codigo}[Andando]
\begin{pilha}
\bk[scEvents]{quando a tecla \menu{seta para cima} for pressionada}
\bk[scMotion]{mude y em \entrada{5}}
\end{pilha}
\begin{pilha}
\bk[scEvents]{quando a tecla \menu{seta para baixo} for pressionada}
\bk[scMotion]{mude y em \entrada{$-5$}}
\end{pilha}
\begin{pilha}
\bk[scEvents]{quando a tecla \menu{seta direita} for pressionada}
\bk[scMotion]{mude x em \entrada{5}}
\end{pilha}
\begin{pilha}
\bk[scEvents]{quando a tecla \menu{seta esquerda} for pressionada}
\bk[scMotion]{mude x em \entrada{$-5$}}
\end{pilha}
\end{codigo}

\section{Passo 3: paredes e meta}

Um laço \bloco[scControl]{sempre} vigia as cores o tempo todo:

\begin{codigo}[Vigiando as cores]
\renewcommand{\sbw}{340pt}
\begin{pilha}
\bk[scEvents]{quando \bandeira\ for clicado}
\bk[scLooks]{mude o tamanho para \entrada{30} \%}
\bk[scMotion]{vá para x: \entrada{$-200$} y: \entrada{150}}
\sbc[scControl]{sempre}
\sbc[scControl]{se \booleano{tocando na cor \menu{preto}?} então}
\bk[scMotion]{vá para x: \entrada{$-200$} y: \entrada{150}}
\sbcend
\sbc[scControl]{se \booleano{tocando na cor \menu{verde}?} então}
\bk[scLooks]{diga \entrada{Venci!}}
\bk[scControl]{pare \menu{tudo}}
\sbcend
\sbcend
\end{pilha}
\end{codigo}

\begin{cuidado}
Para o Scratch reconhecer a cor, você precisa escolher \textbf{exatamente} a mesma cor que usou no desenho, usando o conta-gotas do bloco \bloco[scSensing]{tocando na cor}. Cuidado com tons parecidos!
\end{cuidado}

\section{Passo 4: o cronômetro}

Para saber quanto tempo você levou, o Scratch tem o sensor \variavel{cronômetro}.

\begin{codigo}[Medindo o tempo]
\begin{pilha}
\bk[scSensing]{zere o cronômetro}
\bk[scLooks]{diga \reporter{junte \entrada{Tempo: } \variavel{cronômetro}} por \entrada{3} segundos}
\end{pilha}
\end{codigo}

Coloque o \bloco[scSensing]{zere o cronômetro} no começo do jogo e o \bloco[scLooks]{diga} na vitória.

\begin{curiosidade}
Os melhores jogadores de labirinto de verdade treinam para decorar o caminho. Tente bater o seu próprio recorde!
\end{curiosidade}

\section{Melhorando o labirinto}

\begin{exercicio}[Fantasma assustador]
Adicione um ator ``fantasma'' que anda de um lado para o outro. Se o gatinho tocar nele, volta ao início.
\end{exercicio}

\begin{solucao}
No fantasma:
\begin{pilha}
\bk[scEvents]{quando \bandeira\ for clicado}
\sbc[scControl]{sempre}
\bk[scMotion]{mova \entrada{3} passos}
\bk[scMotion]{se tocar na borda, volte}
\sbcend
\end{pilha}
No gatinho, dentro do \bloco[scControl]{sempre}:
\begin{pilha}
\sbc[scControl]{se \booleano{tocando em \menu{Fantasma}?} então}
\bk[scMotion]{vá para x: \entrada{$-200$} y: \entrada{150}}
\sbcend
\end{pilha}
\end{solucao}

\begin{exercicio}[Contador de tentativas]
Conte quantas vezes o gatinho bateu na parede. Use a variável \textbf{batidas}.
\end{exercicio}

\begin{solucao}
\begin{pilha}
\bk[scVariables]{defina \menu{batidas} para \entrada{0}}
\sbc[scControl]{se \booleano{tocando na cor \menu{preto}?} então}
\bk[scVariables]{mude \menu{batidas} em \entrada{1}}
\bk[scMotion]{vá para x: \entrada{$-200$} y: \entrada{150}}
\sbcend
\end{pilha}
\end{solucao}

\begin{exercicio}[Segunda fase]
Quando o gatinho chegar na meta, troque para um segundo cenário com outro labirinto.
\end{exercicio}

\begin{solucao}
\begin{pilha}
\sbc[scControl]{se \booleano{tocando na cor \menu{verde}?} então}
\bk[scLooks]{mude o cenário para \menu{labirinto 2}}
\bk[scMotion]{vá para x: \entrada{$-200$} y: \entrada{150}}
\sbcend
\end{pilha}
\end{solucao}

\begin{desafio}[Melhor tempo]
Guarde o melhor tempo numa variável \textbf{recorde}. Ao vencer, se o cronômetro for menor que o recorde, atualize o recorde.
\end{desafio}

\begin{solucao}
Comece com um valor grande, como 999:
\begin{pilha}
\bk[scVariables]{defina \menu{recorde} para \entrada{999}}
\end{pilha}
\begin{pilha}
\sbc[scControl]{se \booleano{\variavel{cronômetro} $<$ \variavel{recorde}} então}
\bk[scVariables]{defina \menu{recorde} para \variavel{cronômetro}}
\sbcend
\end{pilha}
\end{solucao}
